%% =====================================================================
%%  Unit 7 — Do you like coffee?  (Dire ce qu'on aime)
%%  Pistes : 7.01–7.04 dialogues · 7.05–7.06 vocabulaire · 7.07–7.08 grammaire
%%           7.09 prononciation · 7.10–7.13 exercices (textes en annexe)
%% =====================================================================

\newcommand{\otIllusUnVII}{\otIllusDuo{coffee/104/86/15/0.10,tea/190/92/20/0.09,pizza/122/14/13/0.08,
  jazz/150/8/11/0.08,cheese/196/22/12/0.07,music/100/50/12/0.07}{jack}{Do you like coffee?}{lea}{Yes, I do!}}

\unite[Do you like\\coffee?]{7}{Do you like coffee?}{Dire ce qu'on aime}{otIllusUnVII}

\situation{Épisode 7}{Mercredi, pause de 11 heures. Jack et Léa prennent un café en bas du studio et
parlent de ce qu'ils aiment… et de ce qu'Emma n'aime pas. Samedi, Léa passe au café où travaille
Sam : il a une proposition à lui faire.}

\begin{objectifs}
\item dire ce que vous aimez, adorez ou détestez ;
\item poser des questions avec \en{do} et \en{does}, et y répondre : \en{Yes, I do. / No, I don't.} ;
\item mettre un verbe au présent simple à la forme négative : \en{I don't}, \en{she doesn't} ;
\item parler de nourriture, de boissons et de loisirs ;
\item reconnaître les formes faibles de \en{do you} et \en{does she} à l'oral.
\end{objectifs}

\auprogramme
  {Exprimer ses goûts, poser des questions}
  {Présent simple : négation et questions\par\vspace{0.6mm}\en{like} + \en{-ing}}
  {Nourriture, boissons, loisirs : 35 expressions}
  {Les formes faibles : \en{do you}, \en{does she}}
  {Petit-déjeuner et \en{coffee shops}}

\semaine{casque/otNavy/Dialogues/Écouter et lire les deux scènes,
  bulle/otGreenDark/Vocabulaire/Apprendre les 35 expressions,
  blocs/otViolet/Grammaire/do et does,
  sons/otOrangeDark/Sons et culture/Formes faibles et coffee shops,
  crayon/otNavy/Exercices/12 exercices corrigés,
  coche/otGreenDark/Bilan/L'essentiel et le test express}

\clearpage
%% ---------------------------------------------------------------------
\rubrique{ecouter}{Au café}

\begin{dialogue}{1}{Do you like coffee?}{Mercredi, 11 heures, dans le \en{coffee shop} au pied du
  studio. Jack et Léa font une pause.}{7.01}{7.02}
\dit{jack}{\hl{Do you like} coffee, Léa?\note{1}}{Tu aimes le café, Léa ?}
\dit{lea}{\hl{Yes, I do}. I love coffee! And you?}{Oui. J'adore le café ! Et toi ?}
\dit{jack}{I love it too. But Emma \hl{doesn't like} coffee.\note{2}}{J'adore ça aussi. Mais Emma n'aime pas le café.}
\dit{lea}{Really? \hl{What does she drink}?\note{3}}{Vraiment ? Qu'est-ce qu'elle boit ?}
\dit{jack}{Tea, of course! She's British!}{Du thé, bien sûr ! Elle est britannique !}
\dit{lea}{Ha ha! \hl{Does she like} French food?}{Ha ha ! Elle aime la cuisine française ?}
\dit{jack}{\hl{Yes, she does}. She loves cheese and croissants.}{Oui. Elle adore le fromage et les croissants.}
\dit{lea}{And you, Jack? Do you like \hl{cooking}?\note{4}}{Et toi, Jack ? Tu aimes cuisiner ?}
\dit{jack}{\hl{No, I don't}. I \hl{don't cook}. I eat sandwiches!}{Non. Je ne cuisine pas. Je mange des sandwichs !}
\dit{lea}{Every day?}{Tous les jours ?}
\dit{jack}{Every day!}{Tous les jours !}
\end{dialogue}

\begin{comprendre}
\pointnote{1} \en{\textbf{Do you like coffee?}} Pour poser une question au présent simple, on place
  \en{do} devant le sujet. \en{do} ne se traduit pas : c'est un « outil » de question. On répond
  \en{Yes, I do} ou \en{No, I don't}.
\pointnote{2} \en{\textbf{Emma doesn't like coffee.}} À la forme négative, on emploie \en{don't}, et
  \en{doesn't} à la 3\textsuperscript{e} personne. Attention : le verbe reprend alors sa forme de
  base, sans \en{-s} : \en{she doesn't like}.
\pointnote{3} \en{\textbf{What does she drink?}} Avec \en{he}, \en{she}, \en{it}, la question se fait
  avec \en{does}, et le verbe perd son \en{-s} : \en{she drinks} → \en{does she drink?}
\pointnote{4} \en{\textbf{Do you like cooking?}} Après \en{like}, \en{love} et \en{hate}, on met
  souvent le verbe en \en{-ing} : \en{I like cooking} = « j'aime cuisiner ».
\end{comprendre}

\begin{dialogue}{2}{What do you do on weekends?}{Samedi, 10 heures. Léa passe au café où Sam
  travaille le matin.}{7.03}{7.04}
\dit{sam}{Hi, Léa! \hl{What do you do} on weekends?\note{5}}{Salut, Léa ! Qu'est-ce que tu fais le week-end ?}
\dit{lea}{I \hl{like walking} in the park, and I \hl{love reading}.}{J'aime me promener au parc, et j'adore lire.}
\dit{sam}{\hl{Do you play} music?}{Tu joues de la musique ?}
\dit{lea}{No, I don't. But I like \hl{listening to} music.}{Non. Mais j'aime écouter de la musique.}
\dit{sam}{\hl{What kind of} music do you like?\note{6}}{Quel genre de musique tu aimes ?}
\dit{lea}{Jazz and French pop. Do you like jazz?}{Le jazz et la pop française. Tu aimes le jazz ?}
\dit{sam}{I love it! My band plays jazz on Fridays. Do you want to come?\note{7}}{J'adore ! Mon groupe joue du jazz le vendredi. Tu veux venir ?}
\dit{lea}{Yes, I'd love to!}{Oui, avec plaisir !}
\end{dialogue}

\begin{comprendre}
\pointnote{5} \en{\textbf{What do you do?}} Le premier \en{do} est l'outil de la question, le second
  est le verbe « faire ». \en{on weekends} = « le week-end ».
\pointnote{6} \en{\textbf{What kind of…?}} « Quel genre de… ? ». \en{listen to} : on écoute
  \textbf{quelque chose} avec \en{to} : \en{I listen to music}.
\pointnote{7} \en{\textbf{Do you want to come?}} « Tu veux venir ? ». La réponse polie et
  enthousiaste : \en{I'd love to!} (« avec plaisir ! ») : vous la retrouverez à l'unité 15.
\end{comprendre}

%% ---------------------------------------------------------------------
\rubrique{mots}{Les mots de l'unité}

Voici les 35 expressions de l'unité. Les aliments en général s'emploient \textbf{sans article} :
\en{I like cheese} (« j'aime le fromage »), \en{She drinks tea} (« elle boit du thé »).

\begin{vocabulaire}{À manger et à boire}{7.05}
\categorie{Les boissons}
\mot{coffee}{ˈkɑːfi}{le café}
\mot{tea}{tiː}{le thé}
\mot{water}{ˈwɑːtɚ}{l'eau}
\mot{milk}{mɪlk}{le lait}
\mot{juice}{dʒuːs}{le jus (de fruit)}
\categorie{Les aliments}
\mot{bread}{bred}{le pain}
\mot{cheese}{tʃiːz}{le fromage}
\mot{eggs}{eɡz}{les œufs}
\mot{meat}{miːt}{la viande}
\mot{chicken}{ˈtʃɪkɪn}{le poulet}
\mot{fish}{fɪʃ}{le poisson}
\mot{rice}{raɪs}{le riz}
\mot{pasta}{ˈpɑːstə}{les pâtes}
\mot{vegetables}{ˈvedʒtəbəlz}{les légumes}
\mot{fruit}{fruːt}{les fruits}
\mot{chocolate}{ˈtʃɑːklət}{le chocolat}
\mot{a sandwich}{ə ˈsænwɪtʃ}{un sandwich}
\mot{French food}{ˌfrentʃ ˈfuːd}{la cuisine française}
\end{vocabulaire}

\begin{vocabulaire}{Les goûts et les loisirs}{7.06}
\categorie{Aimer ou pas}
\mot{like}{laɪk}{aimer, apprécier}
\mot{love}{lʌv}{adorer, aimer beaucoup}
\mot{hate}{heɪt}{détester}
\mot{I don't like it.}{aɪ ˌdoʊnt ˈlaɪk ɪt}{Je n'aime pas ça.}
\mot{What kind of music?}{ˌwʌt ˈkaɪnd əv ˈmjuːzɪk}{Quel genre de musique ?}
\mot{of course}{əv ˈkɔːrs}{bien sûr}
\categorie{Les loisirs}
\mot{reading}{ˈriːdɪŋ}{la lecture, lire}
\mot{cooking}{ˈkʊkɪŋ}{la cuisine, cuisiner}
\mot{walking}{ˈwɑːkɪŋ}{la marche, se promener}
\mot{dancing}{ˈdænsɪŋ}{la danse, danser}
\mot{traveling}{ˈtrævəlɪŋ}{les voyages, voyager}
\mot{listening to music}{ˈlɪsənɪŋ tə ˈmjuːzɪk}{écouter de la musique}
\mot{watching movies}{ˈwɑːtʃɪŋ ˈmuːviz}{regarder des films}
\mot{playing sports}{ˈpleɪɪŋ ˈspɔːrts}{faire du sport}
\mot{on weekends}{ɑːn ˈwiːkendz}{le week-end}
\mot{Do you want to come?}{də jə ˈwɑːnt tə ˈkʌm}{Tu veux venir ?}
\mot{I'd love to!}{aɪd ˈlʌv tuː}{Avec plaisir !}
\end{vocabulaire}

%% ---------------------------------------------------------------------
\rubrique{grammaire}{Le présent simple : négation et questions}

\otIntertitre[otViolet]{1 · La forme négative : \en{don't} et \en{doesn't}}
\begin{center}\otTableau
\begin{tabular}{@{}>{\enOn}l<{\enOff}@{\hspace{8mm}}>{\enOn\bfseries}l<{\enOff}@{\hspace{8mm}}>{\enOn}l<{\enOff}@{\hspace{8mm}}l@{}}
\otTH{Sujet} & \otTH{Auxiliaire} & \otTH{Verbe} & \otTH{Français}\\\midrule
I, you, we, they        & don't (do not)     & like coffee. & je n'aime pas le café\\
he, she, it, Emma       & doesn't (does not) & like coffee. & elle n'aime pas le café\\\bottomrule
\end{tabular}\end{center}

\otIntertitre[otViolet]{2 · Les questions et les réponses courtes}
\begin{center}\otTableau
\begin{tabular}{@{}>{\enOn\bfseries}l<{\enOff}@{\hspace{3mm}}>{\enOn}l<{\enOff}@{\hspace{3mm}}>{\enOn}l<{\enOff}@{\hspace{10mm}}>{\enOn}l<{\enOff}@{}}
\otTH{Auxiliaire} & \otTH{Sujet} & \otTH{Verbe} & \otTH{Réponses courtes}\\\midrule
Do   & you, we, they  & like jazz?  & Yes, I do. / No, I don't.\\
Does & he, she, it    & like jazz?  & Yes, she does. / No, she doesn't.\\\midrule
\multicolumn{4}{@{}l}{\enb{What \textbf{do} you drink? · Where \textbf{does} she work? · What kind of music \textbf{do} you like?}}\\\bottomrule
\end{tabular}\end{center}

\blocaudio{7.07}{Écoutez et répétez}{Une affirmation, puis la même à la forme négative.}

\begin{regle}[Un seul \en{-s} par phrase]
Quand \en{does} ou \en{doesn't} apparaît, c'est lui qui porte la marque de la 3\textsuperscript{e}
personne : le verbe reste à la base.\par\vspace{1.4mm}
\fauxjuste{She doesn't likes tea.}{She doesn't like tea.}
\fauxjuste{Does he works here?}{Does he work here?}
\end{regle}

\blocaudio{7.08}{Écoutez et répétez}{Une question avec \en{do} ou \en{does}, puis sa réponse courte.}

\otIntertitre[otViolet]{3 · \en{like}, \en{love}, \en{hate} + nom ou verbe en \en{-ing}}
\begin{paire}
\begin{regle}[Avec un nom]
\exen{I like coffee.}{J'aime le café.}
\exen{She loves jazz.}{Elle adore le jazz.}
\exen{I love it!}{J'adore ça !}
\end{regle}
\begin{regle}[Avec un verbe en \en{-ing}]
\exen{I like reading.}{J'aime lire.}
\exen{He hates cooking.}{Il déteste cuisiner.}
\exen{Do you like dancing?}{Tu aimes danser ?}
\end{regle}
\end{paire}

\begin{attention}[Les pièges des francophones]
\fauxjuste{I don't like the coffee.}{I don't like coffee.}
\fauxjuste{You like jazz?}{Do you like jazz?}
\fauxjuste{Yes, I like.}{Yes, I do.}
\fauxjuste{I am not like tea.}{I don't like tea.}
\par\vspace{1mm}{\fontsize{10}{12.6}\selectfont Pour parler d'une chose en général, pas d'article :
\en{I like coffee}, \en{Cheese is good}. On dit \en{the coffee} seulement pour un café précis :
\en{The coffee here is great!}\par}
\end{attention}

%% ---------------------------------------------------------------------
\rubrique[Les formes faibles]{sons}{Les formes faibles : \en{do you}, \en{does she}}

À l'oral, les mots « outils » comme \en{do}, \en{does}, \en{to} ou \en{of} ne sont presque pas
accentués. Leur voyelle s'affaiblit en \phon{ə}, et les mots se fondent les uns dans les autres :
\en{Do you like…?} se dit \phon{dəjə ˈlaɪk}, presque « d'ya like ». Ce sont les \textbf{formes
faibles}. Pour comprendre un anglophone, il faut savoir les reconnaître.

\begin{phrases}{7.09}{Écoutez et répétez ces phrases à vitesse normale. Les mots en gras de la
remarque portent l'accent ; tout le reste glisse.}
\phrase{Do you like coffee?}{\phon{dəjə} · \textbf{like}, \textbf{coffee}}
\phrase{Does she like tea?}{\phon{dəʃi} · \textbf{like}, \textbf{tea}}
\phrase{What do you do?}{\phon{ˈwʌdəjə} · \textbf{what}, \textbf{do}}
\phrase{Where does he work?}{\phon{ˈwerdəzi} · \textbf{where}, \textbf{work}}
\phrase{Do they eat meat?}{\phon{dəðeɪ} · \textbf{eat}, \textbf{meat}}
\phrase{I like listening to music.}{\en{to} \phon{tə} · \textbf{like}, \textbf{listening}, \textbf{music}}
\phrase{Yes, I do.}{à la fin, forme pleine : \phon{duː}}
\phrase{No, she doesn't.}{la négation n'est jamais faible : \phon{ˈdʌzənt}}
\end{phrases}

%% ---------------------------------------------------------------------
\rubrique{culture}{Le petit-déjeuner et les \en{coffee shops}}

\begin{paire}
\begin{culture}[\en{American breakfast}]
Le petit-déjeuner américain peut être salé : œufs (\en{eggs}), bacon, \en{pancakes} au sirop
d'érable, \en{bagels}… En semaine, beaucoup se contentent d'un café et d'un bol de céréales
(\en{cereal}), ou d'un muffin acheté en chemin.
\end{culture}
\begin{culture}[Au \en{coffee shop}]
On commande au comptoir et on précise la taille (\en{small}, \en{medium}, \en{large}). On emporte
souvent son café : \en{to go}. Le serveur écrit parfois votre prénom sur le gobelet… avec une
orthographe surprenante !
\end{culture}
\end{paire}

\begin{culture}[\en{Tea or coffee?}]
Les Américains boivent surtout du café ; les Britanniques, du thé, souvent avec un nuage de lait.
Le célèbre \en{afternoon tea} (thé, petits sandwichs et \en{scones}) reste un plaisir d'occasion,
mais la \en{cup of tea} est un rituel quotidien. Et pour proposer un thé : \en{Would you like a cup of tea?}
\end{culture}

%% ---------------------------------------------------------------------
\rubrique{pratique}{Exercices}

\exercice[ecrire,niveau=1]{À table !}{Écrivez le mot anglais.}
\begin{questions}[2][french]
\item le fromage → \enb{\trou[28mm]{cheese}}
\item l'eau → \enb{\trou[28mm]{water}}
\item le pain → \enb{\trou[28mm]{bread}}
\item le poulet → \enb{\trou[28mm]{chicken}}
\item les œufs → \enb{\trou[28mm]{eggs}}
\item le jus → \enb{\trou[28mm]{juice}}
\item le riz → \enb{\trou[28mm]{rice}}
\item le poisson → \enb{\trou[28mm]{fish}}
\end{questions}

\exercice[ecrire,niveau=1]{À la forme négative}{Récrivez chaque phrase à la forme négative.}
\begin{questions}
\item I like tea. \hfill → \trou[72mm]{I don't like tea.}
\item She drinks coffee. \hfill → \trou[72mm]{She doesn't drink coffee.}
\item We eat meat. \hfill → \trou[72mm]{We don't eat meat.}
\item Jack cooks every day. \hfill → \trou[72mm]{Jack doesn't cook every day.}
\item They play tennis. \hfill → \trou[72mm]{They don't play tennis.}
\item Emma likes coffee. \hfill → \trou[72mm]{Emma doesn't like coffee.}
\end{questions}

\exercice[ecrire,niveau=2]{Posez la question}{Transformez chaque phrase en question.}
\begin{questions}
\item You like jazz. \hfill → \trou[72mm]{Do you like jazz?}
\item She works here. \hfill → \trou[72mm]{Does she work here?}
\item They eat fish. \hfill → \trou[72mm]{Do they eat fish?}
\item Sam plays the guitar. \hfill → \trou[72mm]{Does Sam play the guitar?}
\item Jack cooks. \hfill → \trou[72mm]{Does Jack cook?}
\item You drink tea. \hfill → \trou[72mm]{Do you drink tea?}
\end{questions}

\exercice[ecrire,niveau=1]{Réponses courtes}{Répondez d'après les dialogues, par une réponse courte.}
\begin{questions}[2]
\item Does Léa like coffee? — \trou[36mm]{Yes, she does.}
\item Does Emma like coffee? — \trou[36mm]{No, she doesn't.}
\item Does Jack cook? — \trou[36mm]{No, he doesn't.}
\item Does Emma like French food? — \trou[36mm]{Yes, she does.}
\item Does Léa play music? — \trou[36mm]{No, she doesn't.}
\item Does Sam like jazz? — \trou[36mm]{Yes, he does.}
\end{questions}

\exercice[ecrire,niveau=2]{\en{do}, \en{does}, \en{don't} ou \en{doesn't} ?}{Complétez le dialogue.}
\begin{anglais}
\noindent\begin{tabular}{@{}l@{\hspace{6mm}}l@{}}
\otNomDit{emma} & \trou{Do} you like tea, Léa?\\[1.8mm]
\otNomDit{lea}  & Yes, I \trou{do}. But I \trou{don't} drink it every day.\\[1.8mm]
\otNomDit{emma} & \trou{Does} Jack like tea?\\[1.8mm]
\otNomDit{lea}  & No, he \trou{doesn't}. He only drinks coffee!\\[1.8mm]
\otNomDit{emma} & And what \trou{does} Sam drink?\\[1.8mm]
\otNomDit{lea}  & Coffee too. He \trou{doesn't} like tea.\\[1.8mm]
\otNomDit{emma} & Oh no! I \trou{don't} understand Americans!\\
\end{tabular}
\end{anglais}

\exercice[ecrire,niveau=1]{Avec \en{-ing}}{Complétez avec le verbe entre parenthèses, en \en{-ing}.}
\begin{questions}[2]
\item I love \trou[26mm]{reading}. \textit{(read)}
\item She likes \trou[26mm]{cooking}. \textit{(cook)}
\item Do you like \trou[26mm]{dancing}? \textit{(dance)}
\item We love \trou[26mm]{walking}. \textit{(walk)}
\item He hates \trou[26mm]{watching} TV. \textit{(watch)}
\item I like \trou[26mm]{listening} to jazz. \textit{(listen)}
\end{questions}

\exercice[ecouter,ecrire,niveau=2,piste=7.10]{Les formes faibles}{Écoutez chaque question, dite
deux fois, à vitesse normale. Écrivez les mots qui manquent.}
\begin{questions}
\item \trou[36mm]{Do you} like tea?
\item \trou[36mm]{Does she} work here?
\item \trou[36mm]{What do you} do on weekends?
\item \trou[36mm]{Do they} eat meat?
\item \trou[36mm]{Does he} play the guitar?
\item \trou[36mm]{Where do you} work?
\end{questions}

\exercice[ecouter,niveau=2,piste=7.11]{Qui aime quoi ?}{On interroge Emma, Jack et Sam. Écoutez,
puis écrivez \en{yes} ou \en{no} dans chaque case.}
\begin{center}\otTableau
\begin{tabular}{@{}l@{\hspace{12mm}}c@{\hspace{12mm}}c@{\hspace{12mm}}c@{}}
 & \otTH{Coffee} & \otTH{Tea} & \otTH{Cooking}\\\midrule
{\otSemi Emma} & \enb{\trou[14mm]{no}}  & \enb{\trou[14mm]{yes}} & \enb{\trou[14mm]{yes}}\\[1mm]
{\otSemi Jack} & \enb{\trou[14mm]{yes}} & \enb{\trou[14mm]{no}}  & \enb{\trou[14mm]{no}}\\[1mm]
{\otSemi Sam}  & \enb{\trou[14mm]{yes}} & \enb{\trou[14mm]{yes}} & \enb{\trou[14mm]{no}}\\\bottomrule
\end{tabular}\end{center}

\exercice[ecouter,ecrire,niveau=2,piste=7.12]{Dictée à trous}{Mrs.\ Lopez parle cuisine avec Léa.
Écoutez et complétez.}
\begin{anglais}
\noindent\begin{tabular}{@{}l@{\hspace{6mm}}l@{}}
\otNomDit{lopez} & \trou{Do} you like Mexican food, Léa?\\[1.8mm]
\otNomDit{lea}   & Yes, I \trou{do}! I love tacos.\\[1.8mm]
\otNomDit{lopez} & \trou{Does} your family like spicy food?\\[1.8mm]
\otNomDit{lea}   & No, they \trou{don't}. My mother \trou{doesn't} like it.\\[1.8mm]
\otNomDit{lopez} & \trou{What} do you cook in the evening?\\[1.8mm]
\otNomDit{lea}   & I \trou{cook} pasta, or I eat a sandwich!\\[1.8mm]
\otNomDit{lopez} & Come to my place on Sunday. I \trou{make} great tacos!\\
\end{tabular}
\par\vspace{1mm}{\fontsize{8.6}{10}\selectfont\color{otMuted}\foreignlanguage{french}{\textit{spicy} : épicé · \textit{my place} : chez moi · \textit{make} : faire, préparer}}
\end{anglais}

\exercice[ecrire,niveau=2]{Traduisez}{Traduisez en anglais.}
\begin{questions}[1][french]
\item Tu aimes le café ? — Oui. \textit{(réponse courte)}\lignes{1}
\item Emma n'aime pas le café. Elle boit du thé.\lignes{1}
\item Est-ce que Jack cuisine ? — Non. \textit{(réponse courte)}\lignes{1}
\item J'adore lire et voyager.\lignes{1}
\item Qu'est-ce que tu fais le week-end ?\lignes{1}
\item Nous ne mangeons pas de viande.\lignes{1}
\item Quel genre de musique aimes-tu ?\lignes{1}
\item Elle déteste cuisiner.\lignes{1}
\end{questions}

\exercice[ecrire,niveau=3]{Vos goûts}{En six phrases, dites ce que vous aimez, adorez et détestez :
à manger, à boire, et comme loisirs. Utilisez \en{like}, \en{love}, \en{don't like}, \en{hate}.}
\lignes{6}

\exercice[ecouter,parler,niveau=3,piste=7.13]{À vous de parler !}{Au café, Sam vous pose des questions.
Répondez-lui \textbf{à voix haute} pendant chaque pause.}
\begin{center}\otTableau
\begin{tabular}{@{}>{\otSemi\fontsize{8.6}{10}\selectfont\color{otNavy}}l@{\hspace{3mm}}>{\enOn}p{70mm}<{\enOff}@{\hspace{6mm}}p{78mm}@{}}
 & \otTH{Sam dit\ldots} & \otTH{Vous\ldots}\\\midrule
1. & Hey there! Do you like coffee? & répondez par une réponse courte.\\
2. & What do you drink in the morning? & dites ce que vous buvez le matin.\\
3. & Do you like cooking? & répondez, et dites ce que vous cuisinez (ou pas).\\
4. & What kind of music do you like? & donnez un ou deux genres de musique.\\
5. & My band plays on Friday. Do you want to come? & acceptez avec enthousiasme.\\\bottomrule
\end{tabular}\end{center}

%% ---------------------------------------------------------------------
\rubrique{bilan}{Faites le point}

\begin{essentiel}
\begin{minipage}[t]{0.52\linewidth}
\otEssTitre{Les phrases clés}
\renewcommand{\arraystretch}{1.32}
\begin{tabular}{@{}>{\enOn\otSemi\fontsize{9.4}{11}\selectfont}l<{\enOff}@{\hspace{3mm}}>{\fontsize{9.2}{11}\selectfont\color{white!78!otNavy}}l@{}}
Do you like coffee?          & Tu aimes le café ?\\
Yes, I do. / No, I don't.    & Oui. / Non.\\
She doesn't like tea.        & Elle n'aime pas le thé.\\
What does he drink?          & Qu'est-ce qu'il boit ?\\
I love reading.              & J'adore lire.\\
What kind of music?          & Quel genre de musique ?\\
\end{tabular}
\end{minipage}\hfill
\begin{minipage}[t]{0.42\linewidth}
\otEssTitre{Négation}
{\enb{\fontsize{9.4}{12}\selectfont I \textbf{don't} like · she \textbf{doesn't} like}\par}
\otEssTitre{Question}
{\enb{\fontsize{9.4}{12}\selectfont \textbf{Do} you…? · \textbf{Does} she…? + base verbale}\par}
\otEssTitre{Goûts}
{\fontsize{9.2}{11.6}\selectfont \en{like}, \en{love}, \en{hate} + nom ou \en{-ing}\par}
\otEssTitre{Prononciation}
{\fontsize{9.2}{11.6}\selectfont \en{Do you…?} \phon{dəjə} · \en{Yes, I do} \phon{duː}\par}
\end{minipage}
\end{essentiel}

\begin{jesais}
\item dire ce que j'aime, ce que j'adore et ce que je déteste ;
\item poser une question avec \en{do} ou \en{does} et répondre par une réponse courte ;
\item mettre un verbe à la forme négative avec \en{don't} ou \en{doesn't} ;
\item employer \en{like} + verbe en \en{-ing} ;
\item nommer une vingtaine d'aliments et de boissons, et quelques loisirs ;
\item reconnaître \en{do you} et \en{does she} prononcés rapidement.
\end{jesais}

\clearpage
\begin{testexpress}
\qcm[b]{« Tu aimes le thé ? » se dit :}{\en{You like tea?}}{\en{Do you like tea?}}{\en{Like you tea?}}
\qcm[a]{Complétez : \en{She \rule[-0.6mm]{11mm}{0.5pt} like coffee.}}{\en{doesn't}}{\en{don't}}{\en{isn't}}
\qcm[c]{Quelle phrase est correcte ?}{\en{He doesn't likes tea.}}{\en{He don't like tea.}}{\en{He doesn't like tea.}}
\qcm[b]{\en{Do you like jazz?} — }{\en{Yes, I like.}}{\en{Yes, I do.}}{\en{Yes, I am.}}
\qcm[a]{\en{Does Emma drink coffee?} — }{\en{No, she doesn't.}}{\en{No, she don't.}}{\en{No, she isn't.}}
\qcm[c]{« J'aime cuisiner » se dit :}{\en{I like cook.}}{\en{I like the cooking.}}{\en{I like cooking.}}
\qcm[a]{« Je n'aime pas le fromage » (en général) se dit :}{\en{I don't like cheese.}}{\en{I don't like the cheese.}}{\en{I not like cheese.}}
\qcm[b]{Complétez : \en{What \rule[-0.6mm]{9mm}{0.5pt} he do on weekends?}}{\en{do}}{\en{does}}{\en{is}}
\qcm[c]{« le poulet » s'écrit :}{\en{kitchen}}{\en{chiken}}{\en{chicken}}
\qcm[a]{\en{Do you…?} se prononce le plus souvent :}{\phon{dəjə}, très rapide}{\phon{duː juː}, détaché}{\phon{doʊ juː}}
\end{testexpress}

\noindent{\fontsize{9.6}{12}\selectfont\color{otMuted}Niveau visé (CECRL, niveau A1) : « peut
exprimer simplement ses goûts et poser des questions sur ceux des autres ». Les corrigés de l'unité
sont page~\pageref{cor-u7-1} ; les textes des pistes d'exercices, page~\pageref{trans-7.10}.\par}
